% Authored: development workflow.

\chapter{Development Workflow}
\label{ch:development-workflow}

\section{Environment setup}

\begin{lstlisting}[language=bash]
poetry install            # Python package and dev dependencies
cd frontend && npm install   # only for frontend work
\end{lstlisting}

Optional integrations are installed on demand, for example
\texttt{poetry install -E qdrant -E otel}. CI installs no extras, which is what
keeps the default path honest.

\section{Branches}

CI and the dependency-review workflow trigger on pushes and pull requests
targeting \texttt{main} and \texttt{develop}; the security workflow triggers on
\texttt{main} only. The repository therefore has two long-lived branches, with
\texttt{main} carrying the stricter automated scrutiny.

\section{The local loop}

\srcfile{CONTRIBUTING.md} and \srcfile{AGENTS.md} both specify the same
pre-pull-request sequence, and it matches the CI job step for step:

\begin{lstlisting}[language=bash]
poetry run ruff check .
poetry run ruff format --check .
poetry run mypy src
poetry run pytest --cov=feedback_intelligence_agent --cov-fail-under=63
poetry run python scripts/run_demo.py
\end{lstlisting}

\texttt{make ci} runs the same set plus \texttt{poetry build}. For frontend
changes the additional commands are \texttt{npm audit} and
\texttt{npm run build}.

Documentation work adds one command:

\begin{lstlisting}[language=bash]
make docs
\end{lstlisting}

\section{Conventions contributors are asked to follow}

Both instruction files state the same expectations, and they are visible in the
code:

\begin{itemize}[nosep]
  \item keep the default path working without external APIs or private data;
  \item keep modules small, typed, and explicit; do not remove annotations;
  \item add tests for every new behaviour;
  \item use Pydantic models for external request and response schemas;
  \item use \texttt{pathlib.Path} for filesystem paths and
    \texttt{numpy.typing} for arrays;
  \item raise specific exceptions for data-quality problems;
  \item prefer explicit interfaces over hidden global state;
  \item do not introduce heavyweight frameworks without a clear improvement;
  \item never commit secrets, real customer data, generated datasets, or local
    artifacts;
  \item keep each pull request focused on one change.
\end{itemize}

\section{Repository governance}

The \srcfile{.github/} directory carries a \texttt{CODEOWNERS} file, a pull
request template, three issue templates (bug report, feature request,
documentation) with a template chooser configuration, a Dependabot
configuration, and a release-notes categorisation file. Community health files
--- \srcfile{CODE\_OF\_CONDUCT.md}, \srcfile{CONTRIBUTING.md},
\srcfile{SECURITY.md}, \srcfile{SUPPORT.md} --- are present at the root.

Dependabot opens grouped weekly pull requests for three ecosystems: pip
(with runtime and development groups), npm for the frontend, and GitHub
Actions, each limited to five open pull requests.

The repository contains no commit-message convention file, no pre-commit
configuration, and no branch-protection descriptor, so those aspects of the
workflow cannot be documented from repository evidence.

\noevidence
